In this work, we shall be interested in the case of $\beta = 4$, the so-called symplectic case, since the probability measure is invariant under conjugation by elements of the compact symplectic group, and since we prefer to think of an $n\times n$ quaternion matrix as a $2n \times 2n$ matrix where the quaternions are represented by $2 \times 2$ blocks. For such a $2n \times 2n$ matrix, there are, with probability~1, only $n$ distinct eigenvalues since every eigenvalue occurs with multiplicity $2$ (known in the physics literature as Kramers' degeneracy). These $n$ eigenvalues have the joint PDF
\begin{align}\label{pointprocess}
\rho_n(x_1, \dots, x_n) = \frac{1}{\tilde{Z}_n} \prod_{1 \leq i < j \leq n}|x_i - x_j|^4 \prod_{k=1}^{n} {\rm e}^{-2V(x_k)},
\end{align}
where $\tilde{Z}_n $ is a normalisation constant. We shall furthermore restrict to the case of $V$ being a~real polynomial of even degree $D$ and leading coefficient $\gamma > 0$,
\[
V(x) = \gamma x^D + \mathcal{O}\big(x^{D-1}\big),\qquad x \to \infty.
\]
Here the classical approach (see \cite[Chapter 5]{mehta}, \cite[Chapter 6]{loggases}) involves finding an analogue of orthogonal polynomials and this is supplied by \textit{skew-orthogonal polynomials}. Here the inner product $\langle \cdot, \cdot \rangle_2$ is replaced by a skew-symmetric nondegenerate bilinear form $\langle \cdot, \cdot \rangle_4$, where
\begin{equation}\label{skewprod}
\langle P, Q \rangle_4 \overset{\mathrm{def}}{=} \frac{1}{2} \int_{\mathbb{R}} ( P(x) Q^\prime(x)-P^\prime(x) Q(x) ) {\rm e}^{-2V(x)} {\rm d}x
\end{equation}
for $P$ and $Q$ real polynomials. By \enquote{non-degenerate}, we mean the following. Let $\mathscr{P}_k^\mathbb{R}$ be the space of real polynomials of degree $\leq k$. Then we say $\langle \cdot, \cdot \rangle_4$ is non-degenerate if for all $n \in \mathbb{N}$, we have the implication
\begin{align*}
\forall P \in \mathscr{P}_{2n+1}^\mathbb{R} \big( \langle P, Q \rangle_4 = 0,\, \forall Q \in \mathscr{P}_{2n+1}^\mathbb{R} \Longrightarrow P = 0\big).
\end{align*}

Then a sequence of monic polynomials $\{ P_n \}_{n \in \mathbb{N}}$, $P_n(x) = x^n + \mathcal{O}\big(x^{n-1}\big)$, is said to be skew-orthogonal with respect to $\langle \cdot, \cdot \rangle_4$ if the following holds:
\begin{enumerate}\itemsep=0pt
\item[(1)] $\langle P_{2n}, P_{2m} \rangle_4 = \langle P_{2n+1}, P_{2m+1} \rangle_4 = 0$ for all $n,m \in \mathbb{N}$.
\item[(2)] $\langle P_{2n}, P_{2m+1} \rangle_4 = 0$ for all $n, m \in \mathbb{N}$, $n\neq m$.
\end{enumerate}
The existence of such a sequence is guaranteed by non-degeneracy (it can be constructed by a~procedure analogous to the Gram--Schmidt method), however these properties do not uniquely define the sequence because one can replace $P_{2n+1} \mapsto P_{2n+1} + \alpha P_{2n}$ for any constant $\alpha \in \mathbb{R}$. This degree of freedom may be fixed by demanding that the next-to-leading coefficient of $P_{2n+1}$ is equal to a specified value (it is common to choose this value to be zero however this will not be the most convenient choice for us). It can then be verified that with this extra constraint the sequence is unique.

The importance of skew-orthogonal polynomials for random matrix theory can be seen by the following. Let $h_k = \langle P_{2k}, P_{2k+1} \rangle_4$ be the \enquote{skew norm}, where by the non-degeneracy of~${\langle \cdot, \cdot \rangle_4}$ we necessarily have $h_k \neq 0$. Then (see \cite[Propositions~6.1.6 and 6.1.7]{loggases}) all eigenvalue correlation functions associated to the point process~\eqref{pointprocess} may be expressed in terms of the so-called \textit{pre-kernel},
\begin{equation}\label{prekernel}
S_n(x,y) \overset{\mathrm{def}}{=} \sum_{k=0}^{n-1} \frac{ P_{2k}(x) {\rm e}^{-V(x)} \frac{\rm d}{{\rm d}y}\big( P_{2k+1}(y) {\rm e}^{-V(y)} \big) - P_{2k+1}(x) {\rm e}^{-V(x)} \frac{\rm d}{{\rm d}y}\big( P_{2k}(y) {\rm e}^{-V(y)} \big)}{2 h_k}.
\end{equation}
Precisely, from $S_n$ one constructs a $2 \times 2$ matrix \enquote{kernel}, and then the $k$-point correlation function of the ensemble is written as a $2k \times 2k$ Pfaffian, or equivalently, quaternion determinant (see~\cite[Chapter 6]{loggases}). In particular, the 1-point correlation function is given by~\smash{$\rho_n^{(1)}(x) = S_n(x,x)$}. A~similar treatment in terms of skew-orthogonal polynomials, though with a different skew-inner product than \eqref{skewprod}, can be done in the case of real symmetric matrices ($\beta =1$).
